%% parts/p06/forcing-vocabulary.tex --- shared forcing vocabulary before use

\section{Forcing-order vocabulary}
\label{sec:p06-forcing-vocabulary}

The exact source convention at
Definition~\ref{def:p06-forcing-order-orientation} fixes the
stronger-is-smaller orientation, and the exact source definition at
Definition~\ref{def:p06-poset-filter} controls poset filters, dense sets, and
genericity.  This lookup does not create second identities for them.  It adds
only the cardinal notation needed before Chapter~6.5, compatibility/ccc, and
the many-real Cohen poset used by the later supplements
\cite[pp.~201--224]{Jech2003}.

\begin{hbconvention}{def:p06-forcing-cardinal-notation}{Cardinal notation used before Chapter 6.5}
The target architecture places forcing before the full cardinal-arithmetic
entry.  To prevent a forward use from being undefined, only the notation
needed in Chapters~6.3--6.4 is fixed here.  For cardinals
\(\kappa,\lambda\),
\[
  \kappa^\lambda\defeq\abs{{}^{\lambda}\kappa},
  \qquad
  \lambda^{<\kappa}\defeq
    \sup_{\mu<\kappa}\lambda^\mu,
\]
where \({}^{\lambda}\kappa\) is the set of functions from \(\lambda\) to
\(\kappa\), and \(\kappa^+\) is the least cardinal strictly greater than
\(\kappa\).  The alephs needed below are fixed by
\[
  \aleph_0\defeq\omega,
  \qquad
  \aleph_{\alpha+1}\defeq(\aleph_\alpha)^+,
  \qquad
  \aleph_\delta\defeq\sup_{\beta<\delta}\aleph_\beta
  \quad(\delta\text{ a nonzero limit ordinal}).
\]
Thus
\[
  \mathsf{CH}\defeq2^{\aleph_0}=\aleph_1,
  \qquad
  \mathsf{GCH}\defeq
    \forall\alpha\,\bigl(2^{\aleph_\alpha}=\aleph_{\alpha+1}\bigr).
\]
For an ordinal \(\alpha\), a set \(C\subseteq\alpha\) is
\emph{cofinal} iff every \(\beta<\alpha\) is at most some member of \(C\),
and \(\operatorname{cf}(\alpha)\) is the least order type of such a set.
An infinite cardinal \(\kappa\) is \emph{regular} iff
\(\operatorname{cf}(\kappa)=\kappa\).  These are imported lookup clauses,
not theorem or status claims; Definition~\ref{def:p06-cardinal-arithmetic}
and Definition~\ref{def:p06-cofinality-regular-singular} give the full
Chapter~6.5 entries.
\end{hbconvention}

\begin{hbdefinition}{def:p06-ccc}{Compatibility, antichains, and ccc}
Conditions \(p,q\in\mathbb P\) are \emph{compatible} iff some
\(p^\ast\in\mathbb P\) satisfies \(p^\ast\leq p\) and \(p^\ast\leq q\).  An
\emph{antichain} is a subset of pairwise incompatible conditions.  The poset
has the \emph{countable chain condition} (ccc) iff every antichain is
countable.  These are metalinguistic properties of a poset, not filter clauses;
the distinction is used by the preservation and name-counting theorems.
\index{countable chain condition}\index{antichain!forcing}
\end{hbdefinition}

\begin{hbdefinition}{def:p06-add-forcing}{Cohen forcing \(\operatorname{Add}(\omega,\kappa)\)}
For a cardinal \(\kappa\),
\[
  \operatorname{Add}(\omega,\kappa)
  \defeq
  \operatorname{Fn}(\kappa\times\omega,2,{<}\omega),
\]
the finite partial functions from \(\kappa\times\omega\) to \(2\), ordered by
reverse inclusion.  This poset adds \(\kappa\) coordinate Cohen reals; it is
not the one-real forcing unless \(\kappa=1\).  The distinction is needed
because adding one Cohen real does not by itself force \(\neg\mathsf{CH}\).
\index{Cohen forcing}\index{Add forcing@\(\operatorname{Add}(\omega,\kappa)\)}
\end{hbdefinition}
